% Main report.
%
% Built by "make report", which regenerates every figure and table from
% experiments/results/summary.csv first, then runs latexmk, then runs the dash
% check over the compiled PDF. Nothing in the results chapter is typed by hand.
%
% Ground rule 1 applies to this file: never type a double hyphen in prose, it
% typesets as an en dash. Write "1 to 20".

\documentclass[11pt,a4paper]{report}

\usepackage[T1]{fontenc}
\usepackage[utf8]{inputenc}
\usepackage[english]{babel}
\usepackage{lmodern}
\usepackage{microtype}
\usepackage{geometry}
\geometry{margin=2.6cm}

\usepackage{amsmath}
\usepackage{amssymb}
\usepackage{booktabs}
\usepackage{graphicx}
\usepackage{longtable}
\usepackage{caption}
\usepackage{listings}
\usepackage{xcolor}
\usepackage[hidelinks]{hyperref}

\captionsetup{font=small,labelfont=bf}
\setlength{\parskip}{0.4em}
\setlength{\parindent}{0pt}

\definecolor{codebg}{HTML}{F5F5F2}
\lstset{
  basicstyle=\ttfamily\footnotesize,
  backgroundcolor=\color{codebg},
  frame=none,
  breaklines=true,
  showstringspaces=false,
  captionpos=b
}

\graphicspath{{figures/}}

\title{\textbf{Parallel Numerical Lab}\\[0.4em]
  \large One numerical library, seven execution backends,\\
  and an honest account of what each one costs}
\author{Olajide Badejo}
\date{August 2026}

\begin{document}

\maketitle

\begin{abstract}
\noindent
This report describes a C++23 numerical library in which twelve iterative
solvers are written once, against a single problem interface, and run unchanged
over seven execution backends: serial, OpenMP, POSIX threads, a
\texttt{std::jthread} pool, MPI, hybrid MPI with threads, and CUDA. The design
makes identical numerics across backends a tested invariant rather than an
assumption, and in its strong form: with the deterministic reduction mode every
shared memory backend at every worker count produces bit identical iterates, and
the GPU sweeps are bit identical to the CPU too.

That invariant is what makes the measurements comparable. Because every backend
computes the same values, the differences between them are attributable to the
execution model and to nothing else.

The convergence behaviour of the solver zoo is checked against closed form
theory rather than asserted: the Jacobi spectral radius against
$\cos(\pi h)$, the ratio of two between Jacobi and Gauss Seidel iteration
counts, Young's optimal relaxation factor as a genuine minimum, and the change
of growth order from $O(n^2)$ to $O(n)$ that optimal relaxation buys.

The CPU versus GPU comparison is constructed so that the hardware does not
decide the argument. Iteration counts, which are hardware free, are reported
first. Efficiency against each device's own measured STREAM bandwidth is
reported second. Absolute time is reported last and labelled as a property of
one particular pair of devices. The resulting speedup is then decomposed into
the part that is the memory system and the part that is the implementation.
\end{abstract}

\tableofcontents

\chapter{Introduction}

\section{The question}

The question worth asking about parallel programming models is not whether
OpenMP is fast. It is what each model costs on identical work. That question is
hard to answer honestly because the usual comparison changes two things at once:
a study that reimplements an algorithm per model ends up measuring the care that
went into each reimplementation, and one that compares different algorithms on
different devices ends up measuring the algorithms.

This project answers it by construction. One numerical library, twelve iterative
solvers written exactly once against a single problem interface, and seven
execution backends behind a single parallelism interface. A solver does not know
which backend it is running on. A backend does not know which solver it is
running. What varies between measurements is therefore the execution model, and
that is enforced by a test rather than promised in a paragraph.

\section{The invariant that makes it work}

The design commits to something stronger than the usual claim of numerical
agreement. With the deterministic reduction mode:

\begin{itemize}
  \item every shared memory backend, at every worker count, produces
        \emph{bit identical} iterates, not merely close ones;
  \item the GPU sweeps produce bit identical values to the CPU as well;
  \item across MPI rank counts, results agree to reduction tolerance, and the
        reason they cannot be bit identical is documented rather than hidden.
\end{itemize}

This is achievable because floating point addition is not associative and the
library therefore refuses to let the association vary. The chunk grid used by
every reduction is fixed by the problem size alone and never by the worker
count, so the partial sums are always combined in the same order regardless of
which worker produced which one or when. Chapter \ref{ch:parallelisation}
explains the mechanism; the equivalence suite asserts exact equality across
twelve solvers, two problem families, four backends and seven worker counts.

Getting there required disabling fused multiply add contraction on both the host
and the device. The SOR update has the shape $ab + cd$, which a compiler is free
to fuse, and the host was fusing it while the device was not. A claim of
identical numerics that holds only when the compiler happens not to contract is
a much weaker claim than it appears, so contraction is off on both sides. On
kernels that are bandwidth bound this costs nothing measurable.

\section{The solver zoo, and why it is a family}

Nine methods are required by the objectives and three more earn their place.
Eleven of the twelve are one family: every classical stationary method is a
choice of matrix splitting $A = M - N$, and the differences between Richardson,
Jacobi, the Gauss Seidel variants, SOR and the block methods are entirely
differences in the choice of $M$. Chapter \ref{ch:background} presents them that
way, derives each in a few lines from the common framework, and cites the source
for each convergence result.

Conjugate gradient is the twelfth and is deliberately not in that list. It is
not stationary, its optimality argument is different in kind, and presenting it
as the contrast is what makes the family visible as a family.

Every convergence claim in this report is checked as a measurement. The Jacobi
spectral radius is compared against $\cos(\pi h)$; the ratio of Jacobi to
Gauss Seidel iteration counts is compared against the factor of two that Young's
theory predicts; the closed form optimal relaxation factor is tested as a
minimum by perturbing it in both directions; and the change of growth order from
$O(n^2)$ to $O(n)$ is demonstrated by measuring how the counts grow as the grid
doubles.

\section{A comparison the hardware does not decide}

A five point stencil sweep performs three memory operations and about six
floating point operations per unknown. It is bandwidth bound on every machine
ever built. A ratio of runtimes between a CPU and a GPU on such a kernel is
therefore, to within a small correction, a ratio of memory bandwidths wearing
the costume of a statement about parallel programming.

Chapter \ref{ch:methodology} sets out the protocol that avoids this. Iteration
counts come first because they are properties of the mathematics and transfer to
any machine. Efficiency against each device's own measured STREAM triad
bandwidth comes second, because it is dimensionless and says how well a device
is being used rather than which device is larger. Absolute time comes last,
labelled as a property of this specific pair of devices. The measured speedup is
then decomposed into the factor that is the memory system and the factor that is
the implementation, so a reader can see which is which.

Like is compared against like throughout. Natural ordering Gauss Seidel cannot
run on a GPU at all, because it is sequentially dependent; the comparison
therefore pairs GPU Jacobi against CPU Jacobi and GPU red black Gauss Seidel
against CPU red black Gauss Seidel, and reports the cost of the red black
reordering separately so that it is charged to both sides equally.

\section{What the environment allowed, and what it did not}

The target machine has an Intel Core i7-14700K with eight performance cores and
twelve efficiency cores, and everything runs inside WSL2. The guest does not
pass the heterogeneity through: it reports fourteen uniform cores, identical
cache sizes on every logical processor, and no hybrid feature flag. Thread
affinity binds to a guest virtual processor that the hypervisor may place on any
host core.

Rather than assume a mapping that cannot be verified from inside the guest, the
library measures. A probe times an identical kernel on each logical processor
and reports a two group split only when the throughputs separate by a margin
that dominates the noise. The core knee that the objectives ask for is taken
from the aggregate scaling curve, which depends only on how many cores are
engaged and not on which, and therefore survives virtualisation. Section
\ref{sec:knee} reports both what could and what could not be established.

\section{How to read this report}

Chapter \ref{ch:background} derives the solver family and states the convergence
theory. Chapter \ref{ch:parallelisation} describes the backend interface, what
each implementation does underneath, and the communication analysis for the
distributed backends. Chapter \ref{ch:methodology} is the comparison protocol,
written to be difficult to misquote. Chapter \ref{ch:results} contains the
measurements. Chapter \ref{ch:discussion} interprets them, including the parts
that did not come out as expected. Chapter \ref{ch:conclusion} summarises what
the work established and what it did not.

Every number in Chapter \ref{ch:results} is generated from
\texttt{experiments/results/summary.csv} by a script, and every row of that file
carries the backend, worker count, pinning policy, seed and commit hash that
produced it. A number that has not been measured reads \texttt{pending}.

\chapter{Mathematical background}
\label{ch:background}

\section{The splitting framework}

Let $A$ be nonsingular and write $A = M - N$ with $M$ invertible and cheap to
invert. Then $Ax = b$ is equivalent to $Mx = Nx + b$, and the associated
iteration is

\begin{equation}
  x_{k+1} = M^{-1}\left(N x_k + b\right) = x_k + M^{-1} r_k,
  \qquad r_k = b - A x_k .
  \label{eq:splitting}
\end{equation}

Subtracting the exact solution gives the error recurrence
$e_{k+1} = G e_k$ with iteration matrix

\begin{equation}
  G = M^{-1} N = I - M^{-1} A .
  \label{eq:iteration-matrix}
\end{equation}

The iteration converges for every starting vector if and only if
$\rho(G) < 1$, and the asymptotic rate is $\rho(G)$: the error is reduced by
that factor per iteration once the dominant mode dominates
\cite[Ch.~4]{saad2003}. Everything in this chapter, apart from conjugate
gradient, is a choice of $M$.

Throughout, $A = D + L + U$ splits $A$ into its diagonal, strictly lower and
strictly upper triangular parts.

\section{The model problem}

The test operator is the five point discretisation of $-\Delta u = f$ on the
unit square with homogeneous Dirichlet data. On a uniform mesh with spacing
$h = 1/(n+1)$, multiplying through by $h^2$ gives, for interior indices,

\begin{equation}
  4 u_{i,j} - u_{i-1,j} - u_{i+1,j} - u_{i,j-1} - u_{i,j+1} = h^2 f_{i,j} .
  \label{eq:stencil}
\end{equation}

The resulting matrix is symmetric positive definite with eigenvalues
$4 - 2\cos(p \pi h) - 2\cos(q \pi h)$ for $p, q = 1, \dots, n$, and
eigenvectors $\sin(p \pi x)\sin(q \pi y)$. Its extreme eigenvalues are
$O(h^2)$ and close to $8$, so the spectral condition number is $O(h^{-2})$.
This is the canonical model problem precisely because these quantities are known
in closed form, which lets every convergence claim below be checked rather than
believed \cite[Sec.~11.2]{golub2013}.

\paragraph{A trap worth recording.} The obvious manufactured solution
$u = \sin(\pi x)\sin(\pi y)$ is \emph{exactly} the lowest eigenvector of
\eqref{eq:stencil}. With a zero initial guess the residual is then a single
eigenvector, the Krylov subspace conjugate gradient builds is one dimensional,
and the method converges in one iteration at any grid size. That measures the
degeneracy of the problem, not the speed of the method. The same choice is ideal
for the stationary methods, because the slowest decaying mode of the Jacobi
iteration matrix is that same lowest mode. Both right hand sides are therefore
provided and each is used where it is honest: the single mode source for
discretisation error and stationary rates, a seeded spectrally rich source for
anything involving conjugate gradient and for the whole benchmark sweep.

\section{Richardson iteration}

Take $M = I/\omega$. Then \eqref{eq:splitting} collapses to

\begin{equation}
  x_{k+1} = x_k + \omega\left(b - A x_k\right),
\end{equation}

with $G = I - \omega A$. For symmetric positive definite $A$ with eigenvalues in
$[\lambda_{\min}, \lambda_{\max}]$ the eigenvalues of $G$ are
$1 - \omega\lambda$, so the method converges exactly when
$0 < \omega < 2/\lambda_{\max}$, and $\rho(G)$ is minimised at
$\omega = 2/(\lambda_{\min} + \lambda_{\max})$, where it equals
$(\kappa - 1)/(\kappa + 1)$ \cite[Sec.~4.1]{saad2003}.

The library chooses $\omega = 1/\|A\|_{\text{Gershgorin}}$ when none is given.
Gershgorin's bound overestimates $\lambda_{\max}$ by construction, so the
resulting step is always strictly inside the convergence interval. An earlier
implementation estimated $\lambda_{\max}$ by power iteration and diverged,
because the Rayleigh quotient approaches $\lambda_{\max}$ from below and an
unconverged estimate therefore yields a step that is too large.

\section{Jacobi}

Take $M = D$, so $N = -(L+U)$ and componentwise

\begin{equation}
  x_i^{(k+1)} = \frac{1}{a_{ii}}\left(b_i - \sum_{j \neq i} a_{ij} x_j^{(k)}\right).
\end{equation}

Every component of the new iterate depends only on the old one, which is what
makes the sweep embarrassingly parallel and why its result cannot depend on how
the unknowns were partitioned.

On the model problem the Jacobi iteration matrix has spectral radius

\begin{equation}
  \rho_J = \cos(\pi h),
  \label{eq:rho-jacobi}
\end{equation}

so $\rho_J = 1 - O(h^2)$ and the iteration count to a fixed tolerance grows like
$h^{-2}$, that is, like the number of unknowns. Convergence is guaranteed for
any strictly diagonally dominant matrix \cite[Thm.~4.9]{saad2003}.

\section{Gauss Seidel, and the factor of two}

Take $M = D + L$. Then $(D+L)x_{k+1} = b - Ux_k$ is solved by forward
substitution:

\begin{equation}
  x_i^{(k+1)} = \frac{1}{a_{ii}}\left(b_i
    - \sum_{j<i} a_{ij} x_j^{(k+1)}
    - \sum_{j>i} a_{ij} x_j^{(k)}\right).
\end{equation}

The new iterate is used as soon as it is available, which is why the method
converges faster and why the sweep is sequentially dependent.

For a consistently ordered matrix with property A, of which the five point
stencil in natural ordering is the archetype, Young's theorem gives

\begin{equation}
  \rho_{GS} = \rho_J^{\,2}
  \label{eq:rho-gs}
\end{equation}

\cite[Thm.~4.3]{young1971}. Since the iteration count to a fixed tolerance is
$\log(\varepsilon)/\log\rho$, this means Gauss Seidel needs asymptotically
\emph{half} as many iterations as Jacobi. Section \ref{sec:convergence-results}
reports the measured ratio.

Taking $M = D + U$ instead gives the backward sweep. On a symmetric matrix the
two iteration matrices are similar, so they converge at identical rates; they
differ in how information propagates across the grid, which is why the composite
of one forward and one backward sweep is worth having. That composite,
symmetric Gauss Seidel, has preconditioner
$M = (D+L)D^{-1}(D+U)$, which is symmetric whenever $A$ is. That symmetry is the
whole point: it is what a conjugate gradient preconditioner requires and what
the unsymmetric single sweep cannot provide.

\section{Successive over relaxation}

Take $M = D/\omega + L$, giving the Gauss Seidel update extrapolated by $\omega$:

\begin{equation}
  x_i^{(k+1)} = (1-\omega) x_i^{(k)}
    + \frac{\omega}{a_{ii}}\left(b_i
      - \sum_{j<i} a_{ij} x_j^{(k+1)}
      - \sum_{j>i} a_{ij} x_j^{(k)}\right).
\end{equation}

Kahan's theorem gives the necessary condition $0 < \omega < 2$, and the theorem
of Ostrowski and Reich makes it sufficient for symmetric positive definite $A$
\cite[Sec.~4.2]{saad2003}. For a consistently ordered matrix Young's theory
gives the optimum in closed form:

\begin{equation}
  \omega^{*} = \frac{2}{1 + \sqrt{1 - \rho_J^{\,2}}},
  \qquad
  \rho(G_{\omega^{*}}) = \omega^{*} - 1 .
  \label{eq:omega-opt}
\end{equation}

On the model problem $\rho_J = \cos(\pi h)$, so
$\omega^{*} = 2/(1 + \sin(\pi h))$ and the asymptotic rate becomes
$1 - O(h)$ instead of $1 - O(h^2)$. This is a change in the \emph{order} of the
iteration count, from $O(n^2)$ to $O(n)$ on an $n$ by $n$ grid, and it is the
single largest improvement available inside the stationary family
\cite[Ch.~4 to 6]{young1971}.

\section{Red black ordering}

The graph of the five point stencil is bipartite: colour cell $(i,j)$ red when
$i+j$ is even and black when it is odd, and every red cell has only black
neighbours. Reordering by colour puts $A$ in the form

\begin{equation}
  \begin{pmatrix} D_r & C \\ C^{T} & D_b \end{pmatrix},
\end{equation}

with both diagonal blocks diagonal. A sweep over all red cells is therefore a
Jacobi update within that colour and is fully parallel, and the red and black
half sweeps together are exactly Gauss Seidel in the red black ordering
\cite[Sec.~12.4]{saad2003}. This is the only Gauss Seidel style method a GPU can
run, since the natural ordering recurrence offers no parallelism at all
\cite[Ch.~6]{hager2010}.

The reordering changes the iteration matrix, so it converges at a slightly
different, generally slightly worse, per iteration rate. That penalty is
measured in Section \ref{sec:convergence-results} rather than asserted. The red
black ordering remains consistently ordered in Young's sense, so
\eqref{eq:omega-opt} carries over unchanged, which is why red black SOR is the
standard choice for parallel stencil codes.

\section{Block methods}

Partition the unknowns into contiguous blocks and write $A$ in the corresponding
block form. Taking $M$ to be the block diagonal of $A$ gives block Jacobi, in
which each block solves

\begin{equation}
  A_{pp}\, x_p^{(k+1)} = b_p - \sum_{q \neq p} A_{pq}\, x_q^{(k)}
\end{equation}

independently of the others; taking $M$ to be the block lower triangle gives
block Gauss Seidel, in which blocks are visited in order and read the already
updated ones. Point methods are the special case of one unknown per block.

Each diagonal block is solved \emph{exactly}, by a direct method matched to its
structure: dense LU with partial pivoting for the dense problem, and the Thomas
algorithm for the model problem, whose natural block is a grid line and whose
diagonal block is therefore tridiagonal and diagonally dominant. On the model
problem this is line relaxation, and solving each line exactly rather than each
point cuts the iteration count by roughly a factor of two
\cite[Sec.~4.1.1]{saad2003}.

The block count is a solver parameter and never the worker count. That is what
keeps the iterates of the block methods independent of how many workers are
running, which in turn is what lets the equivalence suite compare them bit for
bit across backends.

\section{Conjugate gradient}

For symmetric positive definite $A$, solving $Ax = b$ is equivalent to
minimising $f(x) = \tfrac{1}{2}x^{T}Ax - b^{T}x$, whose gradient is $Ax - b$, so
the residual is the direction of steepest descent. Steepest descent is slow
because successive steps undo one another; conjugate gradient fixes this by
choosing search directions that are $A$ orthogonal, $p_i^{T} A p_j = 0$ for
$i \neq j$. Under that condition an exact line search along each direction never
needs revisiting, so after $k$ steps the iterate is optimal over the Krylov
subspace $\mathrm{span}\{r_0, Ar_0, \dots, A^{k-1}r_0\}$ and, in exact
arithmetic, the method terminates in at most $n$ steps \cite{shewchuk1994}.

The practical recurrence needs one matrix vector product and two inner products
per iteration:

\begin{align}
  \alpha_k   &= \frac{r_k^{T} r_k}{p_k^{T} A p_k}, &
  x_{k+1}    &= x_k + \alpha_k p_k, &
  r_{k+1}    &= r_k - \alpha_k A p_k, \notag \\
  \beta_k    &= \frac{r_{k+1}^{T} r_{k+1}}{r_k^{T} r_k}, &
  p_{k+1}    &= r_{k+1} + \beta_k p_k. &&
  \label{eq:cg}
\end{align}

The error in the $A$ norm satisfies

\begin{equation}
  \|e_k\|_A \le 2\left(\frac{\sqrt{\kappa}-1}{\sqrt{\kappa}+1}\right)^{k}\|e_0\|_A .
  \label{eq:cg-bound}
\end{equation}

On the model problem $\kappa = O(h^{-2})$, so $\sqrt{\kappa} = O(h^{-1})$ and the
iteration count is $O(n)$: the same order as optimally relaxed SOR, but reached
without needing to know a relaxation factor in advance \cite[Ch.~6]{saad2003}.
That is why it displaced the stationary methods in practice.

The two inner products are global reductions, and they are the reason conjugate
gradient scales worse than Jacobi on a distributed machine: each one is a
synchronisation point that no amount of local work can hide. Section
\ref{sec:mpi-results} measures that cost.

\section{The numerics module}

The solvers rest on a small set of classical routines, each with its
convergence order documented and tested: bisection, Newton and Brent for root
finding \cite[Ch.~2]{burden2016}\cite[Ch.~4]{brent1973}; adaptive Simpson,
Gauss Legendre with nodes computed by Newton iteration on the Legendre
recurrence rather than tabulated, and Romberg extrapolation
\cite[Ch.~4]{burden2016}\cite{davis1984}; classical RK4 and adaptive Dormand
Prince 5(4) \cite{dormand1980}\cite[Sec.~II.4]{hairer1993}; and dense LU with
partial pivoting together with Householder QR
\cite[Sec.~3.2, 5.2]{golub2013}\cite[Lec.~10, 16]{trefethen1997}.

Dormand Prince is preferred to Fehlberg because its coefficients are fitted to
minimise the error of the fifth order formula, so the fifth order result is the
one propagated and the embedded fourth order formula serves only as the error
estimate; Fehlberg optimises the fourth order formula instead and thereby wastes
the better solution. Householder reflections are preferred to Gram Schmidt
because they are unconditionally backward stable, whereas classical Gram Schmidt
loses orthogonality in proportion to the condition number.

\chapter{Parallelisation design}
\label{ch:parallelisation}

\section{One interface}

Every backend implements the same small interface, and it is deliberately the
minimal common denominator of the models it spans: a data parallel map, a
reduction, a barrier, and the few extra operations a distributed memory model
genuinely needs. Anything richer would privilege one model and stop the
comparison being fair. The shape follows the structured parallel patterns
treatment of McCool, Robison and Reinders \cite[Ch.~3, 5]{mccool2012}.

\begin{lstlisting}[caption={The parallelism interface, abbreviated}]
void parallel_for(Index n, const RangeBody& body);
Real reduce(Index n, Real init, const RangeReducer& reducer);
void barrier();

Range local_rows(Index total_rows) const;
void exchange_halo(VectorView grid, Index stride, Index rows);
void gather_rows(VectorView data, Range local);
void run_ordered(const std::function<void()>& work, bool forward, ...);
\end{lstlisting}

\subsection{Why the callbacks are chunk level}

A body receives a \texttt{Range} and loops over it itself, rather than being
called once per element. Two consequences follow, both deliberate.

First, the \texttt{std::function} indirection is paid once per chunk rather than
once per element, so it is $O(\text{workers})$ per sweep and does not
contaminate the timings the study exists to produce. Second, the innermost loop
stays a plain loop over contiguous memory that the compiler can still vectorise,
so each backend measures its own dispatch and synchronisation cost rather than a
penalty the abstraction imposed on all of them equally.

\subsection{What backends may not do}

No backend leaks its model's types through the interface: there is no
\texttt{MPI\_Comm}, no \texttt{omp\_} type and no \texttt{cudaStream\_t} in any
signature. The CUDA path goes further and crosses a C ABI boundary, which it has
to for reasons given in Section \ref{sec:cuda-design}.

\section{Determinism as a design decision}
\label{sec:determinism}

Floating point addition is not associative, so the result of a reduction depends
on the order in which the partial sums are combined. Most parallel libraries
treat this as an unavoidable nuisance and paper over it with a tolerance. This
library treats it as a choice.

Under \texttt{ReductionMode::Deterministic}, the chunk grid used by every
reduction is fixed by the problem size alone:

\begin{lstlisting}[caption={The chunking policy that makes determinism possible}]
inline constexpr Index DETERMINISTIC_CHUNKS = 512;

constexpr Index reduction_chunk_count(Index n) noexcept {
    return n <= 0 ? 0 : std::min(DETERMINISTIC_CHUNKS, n);
}
\end{lstlisting}

Each chunk's partial lands in its own slot and the slots are summed in ascending
index order afterwards. Which worker computed which slot, and in what order they
finished, therefore cannot affect the answer. The consequence is that a
reduction computed on one thread, on twenty threads, or on a persistent pool
returns bit identical results, and the equivalence suite can assert exact
equality rather than hiding a real disagreement under a tolerance.

Two limits are documented rather than concealed. Across MPI rank counts the
results agree to reduction tolerance and not bitwise, because rank boundaries
are chosen for load balance and therefore do not align with the fixed chunk
grid, which regroups the additions; within one rank count an ordered allgather
makes the result reproducible and identical on every rank. On the GPU the
reductions use a tree inside each thread block, which is not an ordered sum and
cannot be made into one by a compiler flag; GPU \emph{sweeps} are bit identical
to the CPU, and anything passing through a reduction is compared to tolerance.

The cost of this choice is measured, not waved away: Section
\ref{sec:reduction-cost} prices the deterministic mode against each model's
native reduction.

\section{Contraction, and a claim that nearly was not true}

Getting the GPU sweeps to agree bit for bit with the CPU required disabling
fused multiply add contraction on both sides: \verb|--fmad=false| for
\texttt{nvcc} and \texttt{-ffp-contract=off} for the host compiler. (Those flags
are set in \verb|\verb| spans throughout, because a double hyphen in ordinary
prose typesets as exactly the character ground rule 1 forbids.)

The reason is instructive. The Jacobi update has the form
$\tfrac{1}{4}(b + \sum \text{neighbours})$, which offers no opportunity to fuse,
and it agreed immediately. The SOR update has the form $ab + cd$, which a
compiler may legally contract into a fused multiply add, and the host was
contracting it while the device was not. The disagreement was one bit, and it
would have been easy to absorb into a tolerance and never think about again.

It was not absorbed, because a claim of identical numerics that holds only when
the compiler happens not to contract is a much weaker claim than it appears, and
one that would quietly break on a compiler upgrade. On kernels that are
bandwidth bound the cost of disabling contraction is not measurable.

\section{What each backend does underneath}

\subsection{Shared memory}

\textbf{OpenMP} is the only backend that does not partition the range itself. It
hands the chunk grid to OpenMP's own worksharing construct, so what is measured
is OpenMP's scheduler and barrier rather than a hand written partitioner wearing
an OpenMP hat. The build asserts the \texttt{\_OPENMP} version macro at compile
time. The reference document cited is OpenMP 6.0 \cite{openmp2024}; GCC 16
reports 202111, which is OpenMP 5.2, and the code is deliberately restricted to
what the compiler actually implements.

\textbf{POSIX threads} uses raw \texttt{pthread\_create}, a mutex and condition
variable for dispatch, a counter and a second condition variable for completion,
and \texttt{pthread\_setaffinity\_np} for binding. The pool is persistent:
creating threads per sweep would measure \texttt{pthread\_create}, which is a
well known number and not the one this study is about. Workers wait on a
generation counter rather than a flag, so a task published while a worker was
still finishing the previous one cannot be missed, which is the lost wakeup race
a plain boolean would have.

\textbf{The jthread pool} uses only what C++23 provides: \texttt{std::jthread}
workers, \texttt{std::stop\_token} shutdown, and a pair of \texttt{std::barrier}
phases entered in strict alternation. Using one barrier for both release and
collect would let a fast worker race into the next task before a slow one had
left the previous. Work stealing is deliberately absent: the loops here are
uniform stencil sweeps over contiguous memory, where a static partition is
already balanced and a stealing deque would add per task atomics that buy
nothing. Section \ref{sec:schedule-cost} provides the measurement behind that
choice rather than the assertion of it.

\subsection{Distributed memory}

The MPI backend uses remainder aware block row decomposition, so block sizes
differ by at most one and no row is dropped. The halo exchange is two
\texttt{MPI\_Sendrecv} calls moving one grid row each; every MPI call goes
through a checking macro, so a failure becomes a diagnosed exception carrying
its call site rather than a silently wrong answer.

Each rank allocates the whole grid and computes only its own band. A production
code would allocate the local slab only, and this one would too if memory were
the constraint; at 4096 squared a vector is 134 MB, which fits comfortably. What
matters for the study is that the communication volume is identical either way,
and the choice keeps the solver code identical across every backend.

The installed OpenMPI 5.0.10 advertises MPI 3.1 through \texttt{MPI\_VERSION}
while implementing many MPI 4.0 entry points \cite{mpi2021}. The build probes for
symbols rather than testing the version macro, since guarding on the macro would
disable functionality that demonstrably works.

The hybrid backend \emph{inherits} from the MPI backend rather than
reimplementing it, so the halo exchange, the ordered pass and the deterministic
reduction are literally the same code. Any difference the sweep measures between
the two is therefore the threading and nothing else.

\subsection{Sequential dependence across ranks}

Natural ordering Gauss Seidel and SOR are sequentially dependent by definition.
Preserving that across a row decomposition means rank $r$ may not start until
rank $r-1$ has finished and handed over its boundary row, which is the classical
pipelined Gauss Seidel. The library implements exactly that, through
\texttt{run\_ordered}.

The alternative would have been to substitute a red black reordering silently
whenever more than one worker was present. That would have produced attractive
scaling curves for a method that does not scale, by changing the method. Doing
it honestly means these solvers are bit identical to the serial backend at every
rank count \emph{and} show essentially no speedup, and the report reports both.

\section{The CUDA boundary}
\label{sec:cuda-design}

CUDA 13.3 refuses host compilers newer than GCC 15, and cannot parse GCC 15's
own libstdc++ headers either, because they use the C++23 \texttt{if consteval}
statement that the frontend inside \texttt{nvcc} does not implement. The project
needs GCC 16 for \texttt{<mdspan>} and for OpenMP 5.2.

The resolution is a C ABI boundary. The \texttt{.cu} files are compiled by
\texttt{nvcc} driving GCC 14; everything else is compiled by GCC 16; and every
CUDA entry point is declared \texttt{extern "C"} taking plain pointers and
scalars. No libstdc++ type crosses, so the two compilers never have to agree on
a C++ ABI, only on the platform C ABI, which they do by definition. This also
satisfies the rule that backend files never leak their model's types.

The device solver is not a \texttt{Backend} implementation, and that is the right
call rather than an omission. The entire point of running on a GPU is that state
stays in device memory for the whole solve; forcing it behind
\texttt{parallel\_for} would mean a host callback per chunk, which would measure
PCIe latency and nothing else. The grid crosses once at the start and the answer
once at the end, and the transfer time is reported separately so that a reader
whose use case is a single small solve can add it back.

Natural ordering Gauss Seidel is absent from the device on purpose. It is
sequentially dependent and has no parallelism to offer a wide machine. That is a
result, not a gap in the implementation.

\section{Thread pinning under a hypervisor}

The library offers five binding policies, but two of them depend on being able
to tell a performance core from an efficiency core, and on this machine that
cannot be done from inside the guest.

The host is an i7-14700K with eight performance cores of two threads each and
twelve efficiency cores. The guest reports fourteen uniform cores of two
threads, identical 2048K L2 on every logical processor, \texttt{cpu\_capacity} of
1024 everywhere, no hybrid flag in \texttt{/proc/cpuinfo}, and no
\texttt{cpufreq} directory at all. Separately, \texttt{sched\_setaffinity} binds
a thread to a guest virtual processor, and the hypervisor remains free to place
that processor on any host core: a binding is a hint about the guest, not a
guarantee about silicon.

The library therefore measures instead of asking. A probe times an identical
compute kernel on each logical processor and reports a two group split only when
the throughputs separate by a margin that dominates the within group spread;
otherwise it says so, and the core type policies decline to bind. The knee that
Objective 3 asks for is taken from the aggregate scaling curve, which depends
only on how many cores are engaged and not on which, and is therefore valid
under virtualisation. Section \ref{sec:knee} reports both outcomes.

\chapter{Comparison methodology}
\label{ch:methodology}

This chapter is written to be difficult to misquote. It states what is measured,
what each measurement licenses, and what it does not.

\section{Why the naive comparison is misleading}

The tempting experiment is to run a solver on the CPU, run it on the GPU, divide
the times and print the ratio. Three things go wrong.

\paragraph{It compares different algorithms.} Natural ordering Gauss Seidel
cannot run on a GPU at all. If the CPU runs Gauss Seidel and the GPU runs
Jacobi, the ratio conflates the device with the fact that Jacobi needs about
twice as many iterations to reach the same tolerance.

\paragraph{It compares different amounts of effort.} A carefully written GPU port
measured against a CPU version compiled without vectorisation or threading
produces an impressive number about nothing.

\paragraph{It hides what actually limits the calculation.} A five point stencil
sweep moves three doubles and performs about six floating point operations per
unknown. It is bandwidth bound on every machine ever built, so a ratio of
runtimes is, to within a small correction, a ratio of memory bandwidths dressed
up as a statement about parallel programming.

\section{The protocol}

\subsection{First: iterations to tolerance}

Iteration counts are properties of the mathematics. They are identical on every
machine that runs them, they are what determine whether one method beats another
in principle, and they are therefore reported before any timing. All methods stop
on the same criterion, a relative residual below $10^{-8}$, evaluated the same
way by the same shared driver, so no method is flattered by a laxer test.

\subsection{Like against like}

GPU Jacobi is compared against CPU Jacobi. GPU red black Gauss Seidel is
compared against CPU red black Gauss Seidel. The cost of the red black
reordering, which is a change of iteration matrix and therefore of iteration
count, is measured separately and charged to both sides equally rather than
being allowed to inflate the device figure.

\subsection{Second: bytes per unknown, counted}

A Jacobi sweep writes one value, reads one right hand side entry and reads the
previous iterate. The four stencil neighbours of consecutive unknowns overlap,
so in the streaming limit each is already resident from a nearby access. That is
three doubles, \textbf{24 bytes per unknown per pass}. The same figure is
returned by the host implementation and set by the device solver, so the number
used in the analysis comes from the code rather than from an estimate.

That count is the conservative one and it is not the only defensible one. A
store to a line the cache does not hold has to fetch that line before it can
modify it, which is read for ownership and is a property of a write allocate
cache rather than of the algorithm. A Jacobi pass writes an output array it
never reads, so under that model the count is short by one double per unknown
and the true figure is \textbf{32 bytes}. Which of the two is right is an open
question, so both are published: every result row carries both counts, and every
table and figure here that shows an achieved bandwidth shows both, labelled
\emph{declared} and \emph{counted}. Neither replaces the other, and quoting one
without the other is quoting half the evidence.

\subsection{Pre registered: the rule that settles the traffic model}
\label{sec:preregistered-traffic}

The question is settled by measurement and the rule that reads that measurement
was fixed before it was taken, in
\texttt{benchmarks/sweep\_matrix.yaml} under
\texttt{preregistered.traffic\_model}. Two host STREAM triads run over the same
arrays, the same sizes, the same worker counts and the same repetitions,
differing in the store instruction and in nothing else: a plain C++ loop, whose
compiler emitted stores are the ones under suspicion, and the same triad written
through \texttt{\_mm256\_stream\_pd} with one \texttt{\_mm\_sfence}, which does
not fetch the line it overwrites. Both report against the same declared 24 bytes
per element, so the ratio of the two is the ratio of the traffic they really
move.

The rule, quoted from the registration:

\begin{quote}
A ratio above 1.20 selects the read for ownership model, 32 bytes per unknown
per pass. A ratio below 1.10 selects the conservative model, 24 bytes. A ratio
from 1.10 to 1.20 inclusive is recorded as unresolved and both models are
carried. The thresholds are fixed here, before any measurement, and are not
revisited afterwards.
\end{quote}

The statistic the rule is applied to was amended on 2026-09-06, before the
publication session and before any measurement the rule governs, because the
statistic as first registered had a run to run spread as wide as the undecided
band it had to fall inside, and because its two halves could be taken at
different worker counts. The thresholds and the sentences below were not
touched. As amended: five repetitions of each probe at every worker count,
alternating the two so a repetition's arms see the same machine; $w^*$ is the
worker count at which the plain triad's median is highest; the statistic is the
median over the five repetitions of the non temporal figure over the plain
figure at $w^*$, each ratio taken within one repetition; and if the interval
from the smallest to the largest of those five ratios contains either threshold
the outcome is unresolved whatever the statistic is, which is ground rule 7
applied to a denominator.

\paragraph{Pre registered, one sentence per outcome.} Each of the following was
written before the measurement, and the report carries whichever one the
measurement selects, with the ratio substituted for the placeholder. They are
quoted here verbatim so that a reader can check the report against what was
promised rather than against what was convenient.

\begin{quote}
\textbf{If the read for ownership model is selected.} The non temporal triad
reached a ratio of \texttt{<ratio>} against the plain triad on the publication
machine, above the 1.20 threshold fixed before the measurement, so the read for
ownership model is selected: a Jacobi pass over the five point stencil moves 32
bytes per unknown and not 24, the counted column is the achieved bandwidth this
report compares against each device's own triad, and the host figures rise by a
third while the device figures do not move, which is why host and device
efficiency converge on the Jacobi row.
\end{quote}

\begin{quote}
\textbf{If the conservative model is selected.} The non temporal triad reached a
ratio of \texttt{<ratio>} against the plain triad on the publication machine,
below the 1.10 threshold fixed before the measurement, so the conservative model
is selected: a Jacobi pass moves 24 bytes per unknown, the declared column
stands as the achieved bandwidth of this report, and Section 4.2's read for
ownership argument is refuted on this machine rather than confirmed. The Jacobi
efficiency gap between host and device is then a real gap and not an artefact of
the denominator, and the only correction this release makes to the work unit is
the pass count of phase A2.
\end{quote}

\begin{quote}
\textbf{If it is unresolved.} The non temporal triad reached a ratio of
\texttt{<ratio>} against the plain triad on the publication machine, between the
1.10 and 1.20 thresholds fixed before the measurement, so the traffic model is
recorded as unresolved: both counts, 24 and 32 bytes per unknown per pass, are
carried in every table and figure of this report, neither is presented on its
own as the achieved bandwidth, and no claim is made that rests on one of them
and would fail under the other.
\end{quote}

Which of the three the measurement selected is stated in Chapter
\ref{ch:results}, from the session manifest, by the generator rather than by a
hand that had already seen the number. The publication session applies the rule
for this release; the assembly triad of release 1.2.0 rebuilds the same
experiment and confirms the outcome rather than gating it.

\subsection{Third: each device's own bandwidth, measured}

Both devices are probed with the same STREAM triad kernel,
$a_i = b_i + q\,c_i$, which moves two reads and one write per element and has no
reuse, so it measures the memory system and nothing else \cite{mccalpin1995}.
The host probe runs over the execution backend rather than on one thread,
because a single core cannot saturate a twenty core part and dividing by a
single core figure would make every parallel result look superlinear.

Manufacturer bandwidth figures are never used. They are a bus width times a
clock, no real kernel reaches them, and dividing by one would make both devices
look equally inefficient while saying nothing about which is being used well.

\subsection{The comparable number}

\begin{equation}
  \text{efficiency} =
  \frac{\text{achieved GiB/s on the sweep}}
       {\text{that device's own measured triad GiB/s}} .
  \label{eq:efficiency}
\end{equation}

This is the figure to compare across devices. It is dimensionless, bounded above
by one, and it transfers: a reader with different hardware can measure their own
triad, apply the same fraction, and predict what they would get. It is the
standard roofline reasoning \cite{williams2009} applied to a kernel whose
arithmetic intensity is fixed and known.

\subsection{Last: absolute time}

Seconds are reported last and labelled as a property of one specific CPU paired
with one specific GPU, under one set of compiler flags, on one day. The measured
speedup is then \emph{decomposed} into the factor attributable to the memory
system and the factor attributable to the implementation, so that a reader can
see how much of it is the graphics card and how much of it is the port.

\section{Measuring at sizes where methods do not converge}

Section 8.3 of the specification fixes the comparison at 1024, 2048 and 4096
squared unknowns. At those sizes the stationary methods cannot be run to
convergence: Jacobi on a 1024 squared grid needs of order four million sweeps,
which is days of compute for a number the closed form already predicts.

The two quantities are therefore measured separately, which is more rigorous
than measuring them together rather than less:

\begin{itemize}
  \item \textbf{Iteration counts} are measured at 63, 127 and 255, where running
        to convergence is affordable, and checked against the closed form rates
        of Chapter \ref{ch:background}. It is that agreement which makes
        extrapolating them legitimate rather than a guess. The methods whose
        counts grow like $n$ rather than $n^2$, namely optimally relaxed SOR, its
        red black form and conjugate gradient, are additionally measured at 511
        and 1023, which demonstrates the change of growth order directly.
  \item \textbf{Per iteration cost} is measured at the large sizes with a fixed,
        stated number of sweeps. This is what the bandwidth analysis needs, and
        it is exactly what is affordable there.
\end{itemize}

Time to solution then follows from the two, and because both are tabulated the
crossover point between any two methods on any two devices can be \emph{derived}
rather than asserted:

\begin{equation}
  t = \frac{\text{iterations}(\text{method}, n)\;\times\;n^2\;\times\;\text{bytes per unknown}}
           {\text{achieved bandwidth}(\text{device}, \text{method}, n)} .
  \label{eq:crossover}
\end{equation}

\section{Experimental hygiene}

\begin{itemize}
  \item Every configuration is run once untimed as a warm up, so page faults and
        first touch placement do not land in the measurement, then repeated five
        times with the median reported and the minimum and maximum recorded.
  \item Every result row carries the problem, solver, backend, worker count,
        rank count, pinning policy, reduction mode, schedule, seed and commit
        hash. A row cannot claim provenance the binary did not have: the commit
        is stamped into the binary at configure time and a dirty working tree is
        recorded as such.
  \item Problems are generated from recorded seeds, so any row can be rebuilt.
  \item Both bandwidth probes run once per sweep session and land in the session
        manifest alongside the toolchain versions.
  \item A configuration that fails, or that does not apply, is recorded as such
        with its message rather than dropped.
  \item Compiled with \texttt{-O3 -march=native} and without fast math. Fast math
        would licence the compiler to reassociate the reductions that the
        equivalence suite exists to check.
\end{itemize}

\section{Comparing these numbers with release 1.0.0's}
\label{sec:across-releases}

Every timing in this report replaces the corresponding timing of release 1.0.0
rather than extending it, and none of the values does.

\paragraph{The values are bit identical.} Release 1.1.0 keeps
\texttt{PNL\_REDUCTION\_ACCUMULATORS} at one, so the reduction is the same
ascending chain it was, and every iterate and every residual this library
computes matches release 1.0.0's bit for bit. That is a supported configuration
and not a leftover. The reduction association changes only when that number
changes, which is a later release, and its changelog entry leads with the escape
hatch rather than burying it.

\paragraph{Every timing moved, for three deliberate reasons.} The publication
compiler changed from an unreleased trunk snapshot to a released one, and a
change of compiler changes every timing. The Jacobi solver no longer copies its
whole state on the calling thread once per iteration, so its rows measure less
work rather than the same work faster. The solver's workspace is allocated
outside the timed region, so a repetition no longer pays first touch page fault
cost inside the measurement. The work unit moved with them: sweeps and passes
are separate quantities now, Richardson is charged for one residual evaluation
rather than two, and the symmetric methods are charged the two sweeps they
perform. A 1.1.0 row is therefore not the same quantity as a 1.0.0 row even
where the seconds are close.

The before and after for every headline number that moved is recorded phase by
phase in \texttt{PROGRESS.md}, against the phase that moved it and beside the
gate output that measured it. It is not copied into this chapter, because a
number typed here stops being true at the next sweep, which is the fault this
chapter exists to prevent. The superseded generation keeps its own summary and
its own manifest under \texttt{experiments/results/archive/}, so either can be
rebuilt.

\section{What this comparison does not establish}

\begin{itemize}
  \item \textbf{Nothing about single solve latency.} Transfer time is reported
        separately and is not folded into the bandwidth figure. For one small
        solve, moving the problem across PCIe can cost more than the solve; for a
        simulation that keeps state resident it is paid once.
  \item \textbf{Nothing about other problem classes.} These conclusions hold for
        bandwidth bound stencil work. A dense factorisation is compute bound and
        would rank the devices differently.
  \item \textbf{Nothing about other hardware.} Both measured bandwidths are
        published so that a reader can substitute their own.
  \item \textbf{Nothing about power, cost or programmer time.} Not measured, so
        not claimed.
  \item \textbf{Nothing about performance from the continuous integration
        badge.} CI runs on a shared virtual machine with no GPU; it checks
        correctness, and every timing in this report comes from the target
        machine.
\end{itemize}

\chapter{Results}
\label{ch:results}

Every table and figure in this chapter is generated from
\texttt{experiments/results/summary.csv} by \texttt{scripts/gen\_report\_assets.py},
and so is every measured figure quoted in the prose around them: each one is a
command the generator writes into \texttt{tables/numbers.tex} and this chapter
asks for by name, so a re measurement rewrites the sentences as well as the
tables. Nothing here is typed by hand, and a quantity that has not been measured
reads \texttt{pending} rather than being estimated.

Release 1.0.0 opened this chapter with the same claim above a table that was
typed by hand, and whose figures reconciled with no generated artifact anywhere
in the repository. That table is gone, the mechanism that let it exist is gone
with it, and the claim is now checked: a number the generator cannot derive is a
build failure rather than a sentence that quietly stops being true.

The \pnlConfigurations{} rows behind this chapter come from one sweep at one
commit. Timings from release 1.0.0 are not comparable with these: the compiler,
the Jacobi algorithm and the timed region all changed, and the work unit changed
with them. The values did not. Release 1.1.0 keeps
\texttt{PNL\_REDUCTION\_ACCUMULATORS} at one, so every iterate and every
residual here is bit identical to the one release 1.0.0 computed, and the before
and after for each headline number is recorded phase by phase in
\texttt{PROGRESS.md} rather than copied into this page, where it would stop being
true at the next sweep.

\section{The machine}

\input{tables/bandwidth}

Both figures are measured on this hardware with the same STREAM triad kernel,
and both are what the efficiency numbers later in this chapter divide by. The
host figure is taken across all worker threads, since a single core cannot
saturate the memory system of a twenty core part.

\section{Iterations to tolerance}
\label{sec:convergence-results}

These are the hardware free numbers, and they come first for that reason. They
are identical on any machine that runs them.

\input{tables/convergence}

\begin{figure}[htbp]
  \centering
  \includegraphics[width=0.86\textwidth]{iteration_counts.pdf}
  \caption{Iterations to a relative residual of $10^{-8}$ on the largest grid
    where every method converges in reasonable time. The horizontal scale is
    logarithmic: the span from Richardson to conjugate gradient is more than two
    orders of magnitude, and the ordering of the methods is the entire argument
    of Chapter \ref{ch:background} made visible.}
  \label{fig:iteration-counts}
\end{figure}

\begin{figure}[htbp]
  \centering
  \includegraphics[width=0.82\textwidth]{convergence_growth.pdf}
  \caption{Iteration count against grid size on logarithmic axes, so a power law
    appears as a straight line and its exponent as the slope. The $O(n^2)$
    methods and the $O(n)$ methods separate visibly, which is the practical
    content of Young's optimal relaxation result and of the conjugate gradient
    error bound \eqref{eq:cg-bound}.}
  \label{fig:convergence-growth}
\end{figure}

\subsection{Agreement with theory}

The convergence suite checks each of the following as a measurement against a
closed form prediction rather than accepting it as prose. All pass.

\begin{itemize}
  \item \textbf{The Jacobi spectral radius.} The measured asymptotic contraction
        factor agrees with $\cos(\pi h)$ from \eqref{eq:rho-jacobi} to within
        $10^{-3}$ relative at both 31 and 63 grid points per side.
  \item \textbf{The factor of two.} The convergence suite asserts that the ratio
        of Jacobi to Gauss Seidel asymptotic contraction factors is two within
        two percent, which is Young's $\rho_{GS} = \rho_J^2$ from
        \eqref{eq:rho-gs}, and it passes.
  \item \textbf{Forward against backward.} The two sweep directions agree on
        iteration count to within one percent, as similarity of their iteration
        matrices on a symmetric operator requires.
  \item \textbf{Young's optimum is a minimum.} Perturbing $\omega$ by $\pm 0.02$
        and $\pm 0.05$ about the closed form \eqref{eq:omega-opt} costs
        iterations in both directions.
  \item \textbf{The change of order.} The suite asserts that doubling the grid
        multiplies the Jacobi count by more than three, consistent with
        $O(n^2)$, and the optimally relaxed SOR count by less than 2.6,
        consistent with $O(n)$. Both hold.
  \item \textbf{Conjugate gradient is $O(n)$.} In the swept counts the
        conjugate gradient iteration count grows from \pnlCgIterationsSmall{} to
        \pnlCgIterationsLarge{} as the grid doubles from \pnlCgGridSmall{} to
        \pnlCgGridLarge{}, a factor of \pnlCgGrowthFactor{}.
  \item \textbf{Line beats point.} Solving each grid line exactly rather than
        each point reduces the count, and line Jacobi stands to line Gauss
        Seidel in the same ratio of two as the point methods do.
  \item \textbf{The red black penalty.} Red black Gauss Seidel takes
        \pnlRedBlackIterations{} iterations against natural ordering's
        \pnlNaturalIterations{} on a \pnlRedBlackGrid{} grid, a penalty of
        \pnlRedBlackPenalty{} percent. This is the price of the parallelism, and
        it is charged to both sides of the device comparison.
\end{itemize}

The discretisation itself is second order: the suite fits the order across three
grid sizes and asserts it is two, and asserts the error constant against the
$\pi^2/12 = 0.8225$ predicted by the leading truncation term. Both pass. The
fitted values themselves are a property of the convergence run rather than of
the sweep, so they are reported by that suite and not quoted here.

\section{Backend cost on identical work}

Because every backend computes bit identical values, the differences below are
attributable to the execution model and to nothing else.

\input{tables/backend_cost}

\input{tables/backend_cost_verdicts}

Release 1.0.0 read that table as "the three thread models land within 6 percent
of one another", which states the finding backwards. The data cannot separate
them: on the shared work the differences between OpenMP, POSIX threads and
\texttt{std::jthread} are frequently smaller than the run to run spread of the
rows they are computed from, and where they are, the verdicts above say so
instead of printing a percentage. The conclusion that survives is the weaker and
more useful one, that at this granularity the choice between the three is a
choice about ergonomics rather than about performance, and it is supported by an
inability to tell them apart rather than by a measured gap.

\begin{figure}[htbp]
  \centering
  \includegraphics[width=0.86\textwidth]{backend_cost.pdf}
  \caption{Time for a fixed sweep count at the largest size in the backend cost
    block, all backends at twenty workers. A single hue is used because this is
    a comparison of magnitudes along one categorical axis: the bar length
    carries the measurement and the label carries the identity.}
  \label{fig:backend-cost}
\end{figure}

\section{Scaling and the core knee}
\label{sec:knee}

\begin{figure}[htbp]
  \centering
  \includegraphics[width=0.86\textwidth]{scaling_speedup.pdf}
  \caption{Speedup against worker count for a Jacobi sweep on a 1023 by 1023
    grid. The muted diagonal is ideal scaling and is reference chrome rather
    than a measured series. Line style and marker shape repeat the distinction
    that colour makes, so the figure survives greyscale printing.}
  \label{fig:scaling}
\end{figure}

\input{tables/knee}

\subsection{What could and could not be established}

Objective 3 asks for the performance core versus efficiency core knee on this
heterogeneous CPU. The honest position is as follows.

The per processor probe \textbf{could not} classify logical processors into
performance and efficiency groups. Inside WSL2 the guest reports fourteen
uniform cores, identical cache sizes and no hybrid flag, and thread affinity
binds to a virtual processor the hypervisor may place anywhere. The probe times
an identical kernel on each logical processor and reports a split only when the
throughputs separate by a margin dominating the within group spread; the verdict
it recorded for this session is stored in the session manifest.

The knee \textbf{could} be established from the aggregate scaling curve, which
depends only on how many cores are engaged and not on which, and therefore
survives virtualisation. It is located as the two segment least squares split
that minimises total residual.

\subsection{The knee is not where the objective expected it}

The fitted knee sits at \pnlKneeRange{} workers, not at the eight the objective
anticipated, and the backends do not all put it in the same place. Beyond it the
fitted slope is slightly negative: adding workers does not merely stop helping,
it costs a little.

\input{tables/knee_verdicts}

The disagreement between the backends is a result rather than a defect of the
fit, and it is reported as one. Each position carries the bootstrap interval of
the refits behind it, the union of those intervals spans \pnlKneeInterval{}
workers, and a claim about where one backend's knee sits relative to another's
is only worth making where their intervals do not overlap.

The bandwidth probe explains the shape independently. Sweeping the STREAM triad
across worker counts on an idle machine gives

\input{tables/bandwidth_scaling}

The memory system saturates by \pnlTriadPeakWorkers{} workers and falls away
thereafter, and it is at its worst at \pnlTriadFloorWorkers{}. The curve is flat
enough across the middle of its range that neighbouring points are frequently
not separable at the precision of the repetitions behind them, which the spread
column of that table says row by row. A Jacobi sweep, whose access pattern is
less regular than a triad, saturates at the same point or slightly sooner, which
is the range the knee table gives.

The conclusion is that for this workload \textbf{the binding constraint is the
memory system, not the core types}. The heterogeneous core count that Objective 3
anticipated would set the knee never becomes the limiting factor, because the
machine runs out of memory bandwidth long before it runs out of fast cores. That
is a more useful finding than the one the objective was written to look for, and
it would have been invisible to a study that measured speedup without measuring
the memory system separately.

\paragraph{A claim withdrawn.} An earlier run of this probe peaked at exactly
eight workers, and eight is the performance core count of this processor, which
looked like indirect evidence of the topology the guest hides. Repeating the
measurement on an idle machine moved the peak, and the curve is flat enough
across the middle of its range that the location of its maximum is not stable
between runs. The coincidence does not survive and the inference is withdrawn.
It is recorded here rather than deleted because a plausible pattern that fails to
reproduce is worth exactly one sentence of warning to the next person who notices
it.

\section{Pinning}

\input{tables/pinning}

\input{tables/pinning_verdicts}

Under a hypervisor a binding request constrains guest placement and not silicon,
so the expected effect is smaller than the same experiment would show on bare
metal. The measurement is reported for what it is, policy by policy, and most of
those comparisons are not separable at this precision.

\section{The cost of determinism}
\label{sec:reduction-cost}

Section \ref{sec:determinism} argued that reproducible reductions are worth
having. This is what they cost.

\input{tables/reduction_cost}

\input{tables/reduction_cost_verdicts}

Conjugate gradient is used because it performs two global reductions per
iteration and is therefore the method most exposed to the difference. On a
solver with one reduction per sweep the overhead is proportionally smaller.

Read the spread column before the effect column. Release 1.0.0 quoted this block
as a cost of between minus 5 and plus 3 percent and offered it as what
reproducibility is worth, on runs whose own repetitions moved further than that.
The effects here are mostly smaller than the spread of the rows they are
computed from, so what this table establishes is an upper bound rather than a
price: the deterministic reduction does not cost enough on these backends for
this measurement to see it.

\section{Scheduling}
\label{sec:schedule-cost}

\input{tables/schedule_cost}

\input{tables/schedule_cost_verdicts}

This is the measurement behind leaving work stealing out of the jthread pool. On
a uniform stencil sweep a static partition is already balanced, and the extra
bookkeeping of a dynamic schedule buys nothing. The evidence is only as strong as
the verdicts above make it: where a backend's dynamic against static difference
is smaller than the spread of those two rows, that backend contributes nothing to
the argument, and the sentence above rests on the backends where it is larger.

\section{Distributed scaling and the communication fraction}
\label{sec:mpi-results}

\begin{figure}[htbp]
  \centering
  \includegraphics[width=0.86\textwidth]{mpi_scaling.pdf}
  \caption{Speedup against rank count by method. Jacobi exchanges only halo rows
    and scales; conjugate gradient additionally performs two global reductions
    per iteration, each a synchronisation point no amount of local work can
    hide; natural ordering Gauss Seidel is sequentially dependent and is
    included to show what that costs when the semantics are preserved rather
    than quietly replaced.}
  \label{fig:mpi-scaling}
\end{figure}

The MPI backend records halo, reduction, barrier and ordered pass time
separately, so the communication fraction is measured rather than inferred from
the gap between observed and ideal speedup.

\section{The CPU versus GPU comparison}
\label{sec:device-results}

Read this section in the order it is presented. The iteration counts of Section
\ref{sec:convergence-results} come first and are hardware free. The efficiency
column below is the comparable one. The seconds column is a property of this
particular pair of devices.

\input{tables/device_comparison}

\begin{figure}[htbp]
  \centering
  \includegraphics[width=0.82\textwidth]{device_efficiency.pdf}
  \caption{Percentage of each device's own measured STREAM triad bandwidth
    achieved, method by method. This is the fair comparison: it is dimensionless
    and says how well each device is used, rather than which device is larger.
    The outlined bars are the counted byte model against each triad recounted
    the same way, so they sit on the solid bar for a method that is one pass to
    one sweep and above it, as an upper bound, for one that is not.}
  \label{fig:device-efficiency}
\end{figure}

\subsection{Decomposing the speedup}

The wall clock ratio between the two devices is not a single fact. It factors
into two, and the factorisation is the point:

\begin{equation}
  \underbrace{\frac{t_{\text{host}}}{t_{\text{device}}}}_{\text{speedup}}
  \;=\;
  \underbrace{\frac{B_{\text{device}}}{B_{\text{host}}}}_{\text{memory systems}}
  \;\times\;
  \underbrace{\frac{\eta_{\text{device}}}{\eta_{\text{host}}}}_{\text{how well each is used}},
  \label{eq:decomposition}
\end{equation}

with $B$ the measured triad bandwidth of each device and $\eta$ the efficiency
of \eqref{eq:efficiency}. The first factor was bought with the graphics card and
would be the same for any bandwidth bound kernel. The second is the only part
attributable to the implementations.

On the measurements in the table above, the great majority of the observed
speedup is the first factor. That is the honest summary: on a bandwidth bound
stencil sweep this GPU's measured triad is \pnlDeviceOverHostTriad{} times the
host's, and the device is used somewhat more efficiently while doing it. It is a
useful thing to know and considerably less exciting than an undecomposed speedup
figure, which is why it is presented this way.

\paragraph{What each column bounds.} The two bandwidth columns of Table
\ref{tab:device-comparison} are not two estimates of one quantity. Each
percentage divides by the triad counted the same way as its own numerator, so
the read for ownership charge cancels in the ratio and a counted efficiency is
the declared efficiency times the row's pass count over its sweep count. The
byte model therefore moves the achieved bandwidth figure and never the
efficiency. Jacobi is one pass to one sweep, so at \pnlDeviceSize{} unknowns it
reads \pnlJacobiHostEfficiency{} percent on the host and
\pnlJacobiDeviceEfficiency{} percent on the device under both models alike.

For a method with more passes than sweeps the counted model charges every pass a
full stencil pass, and that overstates the traffic: a red black sweep is two
passes that each write half the unknowns, and conjugate gradient's six passes are
one stencil product, two inner products that write nothing and three vector
updates that read what they write. The counted column is an upper bound for those
methods and the declared column a lower one, so the pair brackets the traffic
rather than competing to describe it, and conjugate gradient's counted
\pnlCgHostCounted{} percent on the host and \pnlCgDeviceCounted{} percent on the
device are a bound that is loose rather than a kernel that outran its memory
system. A per pass model that charges each pass what it actually moves is the
honest next step. It wants the traffic model settled first and is release 1.2.0
work, so this release brackets the traffic instead of estimating it.

\input{tables/traffic_model}

\subsection{The crossover, derived}

Because iteration counts and per iteration costs are measured separately, the
point at which one method on one device overtakes another can be computed from
\eqref{eq:crossover} rather than declared. Both sides of that expression are
tabulated in this chapter, so a reader can substitute their own hardware by
replacing the bandwidth term with a triad measurement of their own.

\chapter{Discussion}
\label{ch:discussion}

\section{What the equivalence invariant bought}

The decision that shaped everything else was to fix the reduction chunk grid by
problem size alone. It looked at first like a small implementation detail and it
turned out to be the foundation of the whole comparison.

Because every shared memory backend produces bit identical iterates, the timing
differences between them cannot be attributed to any numerical difference: there
is none. That removes an entire class of objection from the results chapter. It
also turned the test suite into something sharper than a regression check.
Asserting exact equality means a test fails on the first bit that differs, which
is a much more sensitive detector than a tolerance that would have absorbed real
faults silently. The fused multiply add discrepancy in Section
\ref{sec:cuda-design} was one bit in the last place, and a tolerance based test
would never have reported it.

The cost is real and is measured in Section \ref{sec:reduction-cost} rather than
excused. It is worth paying here because reproducibility is a deliverable of
this project; a production solver with different priorities might reasonably
choose the native reduction and say so.

\section{The methods that do not parallelise, and why they are still here}

Natural ordering Gauss Seidel, SOR and block Gauss Seidel are sequentially
dependent. There was a real temptation to substitute a red black reordering
whenever more than one worker was present, which would have produced attractive
scaling curves for all of them.

That would have been a lie of a particular kind: not a wrong number, but a right
number for a different algorithm. The library instead preserves exact sequential
semantics on every backend, including across MPI ranks through a pipelined token
chain. The consequence is that these methods appear in the scaling figures as
essentially flat lines, and that is the correct result. The reader learns
something true, which is that a method with a sequential recurrence has nothing
to offer a parallel machine, and that red black ordering is the price one pays
to get parallelism back.

Charging that price explicitly also improved the device comparison. Because the
red black penalty is measured (\pnlRedBlackIterations{} against
\pnlNaturalIterations{} iterations on a \pnlRedBlackGrid{} grid,
\pnlRedBlackPenalty{} percent), it can be applied to both the CPU and the GPU
side, and the device figure is not silently inflated by a change of algorithm.

\section{Where the environment set the limits}

Three findings came from the environment rather than from the mathematics, and
all three changed what could be claimed.

\paragraph{The hybrid core topology is invisible from inside WSL2.} The guest
synthesises a homogeneous machine, and thread affinity binds to a virtual
processor the hypervisor may place anywhere. The response was to measure rather
than assume, and to take the knee from the aggregate scaling curve, which does
not depend on identifying individual cores. That is a weaker claim than the
objective originally envisaged and it is the strongest one the environment
supports.

\paragraph{The CUDA toolchain and the host toolchain cannot be the same.}
CUDA 13.3 rejects any host compiler newer than GCC 15 and cannot parse GCC 15's
own headers either, so no released GCC serves as both. Resolving this through a
C ABI boundary was not a workaround so much as a clarification: the boundary was
already required by the rule that backends do not leak their model's types, and
enforcing it at the ABI level made that rule impossible to violate by accident.

\paragraph{MPI version macros describe conformance, not capability.} OpenMPI
5.0.10 advertises MPI 3.1 while implementing many MPI 4.0 entry points. Guarding
on the version macro would have disabled working functionality; assuming the
functions were present would have broken elsewhere. Probing for symbols is the
only approach that is both correct and portable.

\section{Things that were wrong before they were right}

Four faults are worth naming because each was the kind that produces a plausible
number rather than an obvious failure.

\paragraph{The manufactured solution was an eigenvector.} With
$u = \sin(\pi x)\sin(\pi y)$ the Krylov subspace is one dimensional and
conjugate gradient converges in one iteration at every grid size. That looks
like a spectacular result and is a degenerate problem. It was caught because a
test asserted the iteration count should \emph{grow} with the grid, not merely
that the method converged. A test that had only checked convergence would have
passed forever.

\paragraph{A default that was right for one solver and wrong for another.} The
relaxation factor defaulted to one, which is harmless for SOR, where it means
Gauss Seidel, and catastrophic for Richardson, where on the Poisson operator it
gives an iteration matrix with spectral radius seven. The lesson generalises: a
shared options struct wants defaults that mean "ask the component", not defaults
that happen to suit whichever component was written first.

\paragraph{The resumable sweep was not resumable.} The skip test sat after the
work rather than before it, so a restart re-ran every configuration and skipped
only the write. The counters reported plausible numbers throughout, which is why
it survived until a restart was timed rather than merely observed.

\paragraph{A tolerance below the arithmetic.} Two tests asked for a relative
residual that stationary iterations on this operator cannot reach, because over
relaxation raises their rounding limited floor above it. The failure appeared
only when the solver those tests used changed, which made it look like a
regression in the solver rather than an impossible request. Measuring the floor
turned a confusing symptom into a table, and the tests now ask for a tolerance
above it.

\section{What the traffic model changes, and what it does not}
\label{sec:traffic-correction}

The open question of Section \ref{sec:preregistered-traffic} is whether a Jacobi
pass over the five point stencil moves 24 bytes per unknown or 32, the extra
eight being the output line an ordinary store has to fetch before it can
overwrite it. It is worth being exact about what turns on the answer, because
less turns on it than an earlier draft of this section claimed.

An efficiency is one achieved bandwidth over another, and it says something only
when both sides are counted the same way. The triad this report divides by is a
plain C++ loop, so its own store pays exactly the charge the counted model
charges the kernel, and under the counted model the denominator is recounted
too. The charge then cancels, and a counted efficiency is the declared
efficiency times the row's pass count over its sweep count. The byte model moves
the achieved bandwidth figure in GiB/s and leaves the efficiency where it was,
on the host and on the device alike. Host Jacobi stands at
\pnlJacobiHostEfficiency{} percent of the plain triad and the device at
\pnlJacobiDeviceEfficiency{} percent of its own, under either model, and the two
percentage columns of Table \ref{tab:device-comparison} agree on those rows to
the digit they print.

The Jacobi row is therefore not an exception waiting for a byte model to remove
it, and it is not an exception at all. Its device over host efficiency ratio is
\pnlJacobiEfficiencyRatio{} in this generation, where release 1.0.0's published
generation put it far higher; the before and after per phase is in
\texttt{PROGRESS.md} rather than typed here. What moved it is not the traffic
model. It is that the per iteration full state copy is gone from the Jacobi path
and that the timed region no longer holds the allocation, so the host figure is
a measurement of the sweep rather than of the sweep plus a serial copy plus a
workspace. That is the correction that mattered, and it is a correction to the
measurement and not to the model that reads it.

The other three kernels gave that conclusion independently and Jacobi now joins
them. At \pnlDeviceSize{} unknowns the device over host efficiency ratio is
\pnlRedBlackEfficiencyRatio{} for red black Gauss Seidel and
\pnlCgEfficiencyRatio{} for conjugate gradient, against Jacobi's
\pnlJacobiEfficiencyRatio{}, so on these kernels the great majority of the
observed speedup is the memory system and little of it is attributable to the
port. That holds under the conservative count, without the correction the
traffic model question was about, which is a better outcome than needing it.

\paragraph{The sentence this replaces.} An earlier draft of this section said
that under read for ownership the host achieved bandwidth and the host
efficiency both rise by a third while the device figures do not move, and that
host and device efficiency converge on the Jacobi row as a result. The first
half is right about GiB/s and wrong about the efficiency, and the rest was the
same mismatch: it divided a numerator counted one way by a denominator counted
the other. The pre registered sentence quoted in Section
\ref{sec:preregistered-traffic} carries the same reasoning, and it is quoted
there exactly as it was written, because a pre registration edited after the
measurement is not one. This paragraph is the correction to it, and the outcome
that sentence was written for is not the one the measurement selected.

The report does not choose between the two models on its own authority. The rule
was fixed before the measurement, the measurement is the publication session's,
and the generator applies the rule and writes the sentence, which Section
\ref{sec:device-results} carries. That sentence records the outcome as
\pnlTrafficModel{}, which is this instrument saying it cannot separate the two
counts on this machine rather than choosing between them.

Both counts are therefore carried in every table and every figure, every claim
in this chapter that would change under the other model is stated conditionally,
and no conclusion here rests on one of them in a way that would fail under the
other. For a method whose pass count exceeds its sweep count the two columns
bracket the traffic rather than compete to describe it, which Section
\ref{sec:device-results} sets out; for Jacobi they are the same number.

\section{Related work}
\label{sec:related-work}

Deterministic reduction is established technique and nothing here invented it.
Saying so is not modesty: the contribution of this project is a narrower thing
than the reduction, and not naming the prior art both reads as unawareness of
the field and undersells the work.

\paragraph{Reproducible arithmetic.} Intel's Math Kernel Library offers
Conditional Numerical Reproducibility \cite{mklcnr}, which gives bitwise
identical results across runs and across a range of processors subject to
conditions the vendor documents: a fixed thread count, aligned arrays, and a
selected code path. It states the shape of the problem clearly, which is that
reproducibility is bought by fixing what a performance library would otherwise
be free to vary. Demmel and Nguyen give summation algorithms whose result does
not depend on the order or the number of summands, at bounded cost and with an
error bound \cite{demmel2013,demmel2015}, and Ahrens, Demmel and Nguyen develop
the binning algorithm behind ReproBLAS \cite{ahrens2020}. That is a stronger
guarantee than the one made here: ReproBLAS is reproducible across different
decompositions, whereas the fixed chunk grid of this library is reproducible
because the decomposition is held fixed.

\paragraph{Performance portability.} Kokkos
\cite{edwards2014,trott2022} and RAJA \cite{beckingsale2019} are the established
C++ layers for writing a kernel once and running it over several execution
models including GPUs. Both provide reductions and both document what is and is
not deterministic about them. They solve the portability problem this library
also solves, at a scale and maturity this library does not approach.

\paragraph{Solver libraries.} Ginkgo \cite{anzt2022} is a modern C++ sparse
linear algebra framework built around linear operators, with GPU backends and a
strong benchmarking culture. PETSc \cite{balay1997,petsc2026} is the reference
distributed solver toolkit, whose Krylov methods and preconditioners this
repository implements a small subset of. hypre \cite{falgout2002} is the
reference library for scalable multigrid preconditioners. Each is a real answer
to the question of what to use to solve a system, and this library is not.

\paragraph{What this project adds.} Not the deterministic reduction. Three
things, and they are narrower.

\begin{enumerate}
  \item \textbf{One numerical invariant held across seven execution models at
        once}, a GPU and a distributed backend among them. The comparison
        between programming models in Chapter \ref{ch:results} is only
        meaningful because the thing being executed is provably identical, and
        holding one contract across that particular set of models is the part
        that is unusual.
  \item \textbf{Asserted as exact equality in a test suite}, rather than
        described in documentation. The equivalence suite compares with
        \texttt{==} across twelve solvers, two problem families, four backends
        and seven worker counts, and fails on the first bit that differs. The
        one bit contraction discrepancy of Section \ref{sec:cuda-design} was
        found that way and would have vanished into any tolerance.
  \item \textbf{Priced.} What the invariant costs is measured against each
        model's own native reduction and reported with its run to run spread,
        including where that spread is wide enough that the honest answer is
        that the cost cannot be separated from the noise.
\end{enumerate}

\paragraph{What is missing, and when it arrives.} There is no baseline
comparison in this report. Nothing here says what a competent existing
implementation achieves on this machine on this problem, so a statement of the
form "this solver reaches some percentage of host STREAM" cannot be read as good
or bad. That is a real gap and it is stated as one. A single baseline against
PETSc is planned for release 1.2.0: the same stencil, \texttt{KSPCG} with
\texttt{PCJACOBI} and with \texttt{PCSOR}, the same tolerance, the same machine
and the same reported metrics, emitting rows in the schema the sweep already
uses. One baseline is enough, and the result is worth publishing whichever way
it comes out. \texttt{docs/RELATED\_WORK.md} carries this section in longer form.

\section{What a reader should take away}

For someone choosing a programming model, the results chapter says what each
model costs on identical work, with the numerical variable removed by
construction. For someone choosing a method, the convergence tables say what each
one costs in iterations, independent of any machine. For someone weighing a GPU,
Section \ref{sec:device-results} separates the part of the speedup that is the
memory system from the part that is the implementation, which is the separation
that determines whether the speedup would transfer to their problem.

The most transferable conclusion is methodological rather than numerical. A
comparison between execution models is only meaningful if the thing being
executed is identical, and the only way to know that is to test it at the bit
level. Everything else in this project follows from taking that seriously.

\chapter{Conclusion}
\label{ch:conclusion}

\section{What was built}

A C++23 numerical library in which twelve iterative solvers are written once,
against one problem interface, and run unchanged over seven execution backends.
Around them: a numerics module with root finders, quadrature rules, ODE
integrators and dense factorisations, each with its convergence order tested
against theory; two problem families; a benchmark harness that resolves a
declared matrix, resumes after interruption and records the provenance of every
row; and three reports built from the measurements by script.

\section{What was established}

\begin{enumerate}
  \item \textbf{Identical numerics across execution models is achievable, not
        merely approachable.} With deterministic reductions every shared memory
        backend at every worker count produces bit identical iterates, and the
        GPU sweeps match the CPU bit for bit as well. This is asserted by the
        test suite as exact equality across twelve solvers, two problem
        families, four backends and seven worker counts, and it required fixing
        the reduction chunk grid by problem size and disabling fused multiply
        add contraction on both host and device.

  \item \textbf{The classical convergence theory holds to the digit.} The Jacobi
        spectral radius matches $\cos(\pi h)$; Gauss Seidel needs exactly half
        the iterations of Jacobi, as $\rho_{GS} = \rho_J^2$ requires; Young's
        closed form relaxation factor is a genuine minimum; optimal relaxation
        changes the growth order from $O(n^2)$ to $O(n)$; and conjugate gradient
        reaches the same order without being told a parameter. The five point
        stencil is second order accurate with the error constant $\pi^2/12$
        predicted by the leading truncation term.

  \item \textbf{Sequential methods do not parallelise, and saying so is more
        useful than hiding it.} Natural ordering Gauss Seidel keeps exact
        semantics on every backend, including across MPI ranks, and shows
        essentially no speedup. Red black ordering is what buys the parallelism,
        at a measured cost of 2.6 percent more iterations.

  \item \textbf{A GPU speedup on a bandwidth bound kernel is mostly the memory
        system.} Decomposed through \eqref{eq:decomposition}, the observed
        advantage separates into a large factor that is the difference in
        measured bandwidth and a small factor that is the difference in how well
        each device is used. Reporting the undecomposed ratio would have been
        accurate and misleading.
\end{enumerate}

\section{What was not established}

Stated as plainly as the conclusions, because a limitation left implicit reads
as a claim.

\begin{itemize}
  \item \textbf{The performance versus efficiency core split could not be
        identified from inside the guest.} WSL2 synthesises a homogeneous
        topology and affinity binds to a virtual processor. The knee is reported
        from the aggregate scaling curve instead, which is a weaker claim than
        originally intended and the strongest the environment supports.
  \item \textbf{Nothing is established about compute bound problems.} Every
        conclusion about devices here concerns bandwidth bound stencil work. A
        dense factorisation would rank the hardware differently.
  \item \textbf{Nothing is established about single solve latency}, power, cost,
        or programmer effort. Transfer time is reported separately rather than
        folded in; the rest were not measured and are not claimed.
  \item \textbf{The iteration counts at the largest grids are extrapolated}, not
        measured, for the methods whose counts grow like $n^2$. The
        extrapolation rests on closed form rates that were verified at the sizes
        where measurement was affordable, and the report says which numbers are
        which.
\end{itemize}

\section{Further work}

\paragraph{Preconditioned conjugate gradient.} The symmetric splittings are
already implemented and symmetric Gauss Seidel and SSOR are exactly the
preconditioners the method wants. Combining them would be a small amount of code
and would demonstrate the point of the symmetric variants, which currently exist
in the zoo without the application that motivates them.

\paragraph{Multigrid.} Every component is present: smoothers in the relaxation
sweeps, an exact coarse solve in the LU routine, and the grid structure in the
Poisson problem. Multigrid is the method that makes the iteration count
independent of the grid entirely, and adding it would put the $O(n^2)$ and
$O(n)$ results of this report in their proper context as the two things it
improves on.

\paragraph{A local slab MPI decomposition.} Each rank currently allocates the
whole grid. The communication volume is already correct, so this would not change
any timing reported here, but it would remove the memory ceiling that keeps the
distributed runs below the sizes the device comparison uses.

\paragraph{Bare metal measurement of the core knee.} The one objective the
environment prevented. Running the same sweep outside a hypervisor would settle
whether the knee in the aggregate curve coincides with the performance core
count, which is currently inferred rather than confirmed.

\paragraph{Mixed precision.} The stencil sweeps are bandwidth bound, so halving
the width of the iterate would come close to halving the time. Whether the
resulting convergence penalty is worth it is exactly the kind of question this
library's separation of iteration count from per iteration cost is built to
answer.


\bibliographystyle{plain}
\bibliography{refs}

\end{document}
